% Options for packages loaded elsewhere
\PassOptionsToPackage{unicode}{hyperref}
\PassOptionsToPackage{hyphens}{url}
%
\documentclass[
  11pt,
  letterpaper,
]{article}
\usepackage{amsmath,amssymb}
\usepackage{iftex}
\ifPDFTeX
  \usepackage[T1]{fontenc}
  \usepackage[utf8]{inputenc}
  \usepackage{textcomp} % provide euro and other symbols
\else % if luatex or xetex
  \usepackage{unicode-math} % this also loads fontspec
  \defaultfontfeatures{Scale=MatchLowercase}
  \defaultfontfeatures[\rmfamily]{Ligatures=TeX,Scale=1}
\fi
\usepackage{lmodern}
\ifPDFTeX\else
  % xetex/luatex font selection
  \setmainfont[]{Latin Modern Roman}
  \setmonofont[]{Latin Modern Mono}
\fi
% Use upquote if available, for straight quotes in verbatim environments
\IfFileExists{upquote.sty}{\usepackage{upquote}}{}
\IfFileExists{microtype.sty}{% use microtype if available
  \usepackage[]{microtype}
  \UseMicrotypeSet[protrusion]{basicmath} % disable protrusion for tt fonts
}{}
\makeatletter
\@ifundefined{KOMAClassName}{% if non-KOMA class
  \IfFileExists{parskip.sty}{%
    \usepackage{parskip}
  }{% else
    \setlength{\parindent}{0pt}
    \setlength{\parskip}{6pt plus 2pt minus 1pt}}
}{% if KOMA class
  \KOMAoptions{parskip=half}}
\makeatother
\usepackage{xcolor}
\usepackage[margin=0.8in]{geometry}
\usepackage{color}
\usepackage{fancyvrb}
\newcommand{\VerbBar}{|}
\newcommand{\VERB}{\Verb[commandchars=\\\{\}]}
\DefineVerbatimEnvironment{Highlighting}{Verbatim}{commandchars=\\\{\}}
% Add ',fontsize=\small' for more characters per line
\newenvironment{Shaded}{}{}
\newcommand{\AlertTok}[1]{\textcolor[rgb]{1.00,0.00,0.00}{\textbf{#1}}}
\newcommand{\AnnotationTok}[1]{\textcolor[rgb]{0.38,0.63,0.69}{\textbf{\textit{#1}}}}
\newcommand{\AttributeTok}[1]{\textcolor[rgb]{0.49,0.56,0.16}{#1}}
\newcommand{\BaseNTok}[1]{\textcolor[rgb]{0.25,0.63,0.44}{#1}}
\newcommand{\BuiltInTok}[1]{\textcolor[rgb]{0.00,0.50,0.00}{#1}}
\newcommand{\CharTok}[1]{\textcolor[rgb]{0.25,0.44,0.63}{#1}}
\newcommand{\CommentTok}[1]{\textcolor[rgb]{0.38,0.63,0.69}{\textit{#1}}}
\newcommand{\CommentVarTok}[1]{\textcolor[rgb]{0.38,0.63,0.69}{\textbf{\textit{#1}}}}
\newcommand{\ConstantTok}[1]{\textcolor[rgb]{0.53,0.00,0.00}{#1}}
\newcommand{\ControlFlowTok}[1]{\textcolor[rgb]{0.00,0.44,0.13}{\textbf{#1}}}
\newcommand{\DataTypeTok}[1]{\textcolor[rgb]{0.56,0.13,0.00}{#1}}
\newcommand{\DecValTok}[1]{\textcolor[rgb]{0.25,0.63,0.44}{#1}}
\newcommand{\DocumentationTok}[1]{\textcolor[rgb]{0.73,0.13,0.13}{\textit{#1}}}
\newcommand{\ErrorTok}[1]{\textcolor[rgb]{1.00,0.00,0.00}{\textbf{#1}}}
\newcommand{\ExtensionTok}[1]{#1}
\newcommand{\FloatTok}[1]{\textcolor[rgb]{0.25,0.63,0.44}{#1}}
\newcommand{\FunctionTok}[1]{\textcolor[rgb]{0.02,0.16,0.49}{#1}}
\newcommand{\ImportTok}[1]{\textcolor[rgb]{0.00,0.50,0.00}{\textbf{#1}}}
\newcommand{\InformationTok}[1]{\textcolor[rgb]{0.38,0.63,0.69}{\textbf{\textit{#1}}}}
\newcommand{\KeywordTok}[1]{\textcolor[rgb]{0.00,0.44,0.13}{\textbf{#1}}}
\newcommand{\NormalTok}[1]{#1}
\newcommand{\OperatorTok}[1]{\textcolor[rgb]{0.40,0.40,0.40}{#1}}
\newcommand{\OtherTok}[1]{\textcolor[rgb]{0.00,0.44,0.13}{#1}}
\newcommand{\PreprocessorTok}[1]{\textcolor[rgb]{0.74,0.48,0.00}{#1}}
\newcommand{\RegionMarkerTok}[1]{#1}
\newcommand{\SpecialCharTok}[1]{\textcolor[rgb]{0.25,0.44,0.63}{#1}}
\newcommand{\SpecialStringTok}[1]{\textcolor[rgb]{0.73,0.40,0.53}{#1}}
\newcommand{\StringTok}[1]{\textcolor[rgb]{0.25,0.44,0.63}{#1}}
\newcommand{\VariableTok}[1]{\textcolor[rgb]{0.10,0.09,0.49}{#1}}
\newcommand{\VerbatimStringTok}[1]{\textcolor[rgb]{0.25,0.44,0.63}{#1}}
\newcommand{\WarningTok}[1]{\textcolor[rgb]{0.38,0.63,0.69}{\textbf{\textit{#1}}}}
\usepackage{longtable,booktabs,array}
\usepackage{calc} % for calculating minipage widths
% Correct order of tables after \paragraph or \subparagraph
\usepackage{etoolbox}
\makeatletter
\patchcmd\longtable{\par}{\if@noskipsec\mbox{}\fi\par}{}{}
\makeatother
% Allow footnotes in longtable head/foot
\IfFileExists{footnotehyper.sty}{\usepackage{footnotehyper}}{\usepackage{footnote}}
\makesavenoteenv{longtable}
\setlength{\emergencystretch}{3em} % prevent overfull lines
\providecommand{\tightlist}{%
  \setlength{\itemsep}{0pt}\setlength{\parskip}{0pt}}
\setcounter{secnumdepth}{5}
\ifLuaTeX
\usepackage[bidi=basic]{babel}
\else
\usepackage[bidi=default]{babel}
\fi
\babelprovide[main,import]{american}
\ifPDFTeX
\else
\babelfont{rm}[]{Latin Modern Roman}
\fi
% get rid of language-specific shorthands (see #6817):
\let\LanguageShortHands\languageshorthands
\def\languageshorthands#1{}
\usepackage{amsmath,amssymb}
\usepackage{microtype}
\usepackage{longtable,booktabs,array}
\usepackage{fvextra}
\usepackage{xurl}
\usepackage{fancyhdr}
\usepackage{titlesec}
\usepackage{float}
\AtBeginDocument{\hypersetup{colorlinks=true,linkcolor=black,urlcolor=blue}}
\setlength{\emergencystretch}{3em}
\setlength{\parindent}{0pt}
\setlength{\parskip}{5pt}
\setlength{\headheight}{15pt}
\pagestyle{fancy}
\fancyhf{}
\fancyhead[L]{\small Yang Mills spectral gap foundations}
\fancyhead[R]{\small Review packet}
\fancyfoot[L]{\small 20 November 2026}
\fancyfoot[R]{\small\thepage}
\renewcommand{\headrulewidth}{0pt}
\renewcommand{\footrulewidth}{0pt}
\titleformat{\section}{\Large\bfseries}{\thesection}{0.7em}{}
\titleformat{\subsection}{\large\bfseries}{\thesubsection}{0.7em}{}
\titlespacing*{\section}{0pt}{17pt}{7pt}
\titlespacing*{\subsection}{0pt}{12pt}{5pt}
\let\oldtableofcontents\tableofcontents
\renewcommand{\tableofcontents}{\oldtableofcontents\clearpage}
\DefineVerbatimEnvironment{Highlighting}{Verbatim}{breaklines,commandchars=\\\{\},fontsize=\small}
\widowpenalty=10000
\clubpenalty=10000
\ifLuaTeX
  \usepackage{selnolig}  % disable illegal ligatures
\fi
\usepackage{bookmark}
\IfFileExists{xurl.sty}{\usepackage{xurl}}{} % add URL line breaks if available
\urlstyle{same}
\hypersetup{
  pdftitle={Yang Mills spectral gap foundations},
  pdfauthor={Prepared with ChatGPT for Adam Pippert},
  pdflang={en-US},
  hidelinks,
  pdfcreator={LaTeX via pandoc}}

\title{Yang Mills spectral gap foundations}
\usepackage{etoolbox}
\makeatletter
\providecommand{\subtitle}[1]{% add subtitle to \maketitle
  \apptocmd{\@title}{\par {\large #1 \par}}{}{}
}
\makeatother
\subtitle{Mathematician review packet}
\author{Prepared with ChatGPT for Adam Pippert}
\date{20 November 2026}

\begin{document}
\maketitle

{
\setcounter{tocdepth}{2}
\tableofcontents
}
\section{Review mandate}\label{review-mandate}

This packet asks an independent mathematician to assess a specific
research route: quantitative boundary influence for selected
observables, transfer to temporal clustering, and exclusion of
low-energy spectral weight in a reconstructed quantum theory. The
intended application is four-dimensional pure Yang-Mills theory. The
immediate deliverables are an assessment of the actual hypotheses,
corrections to the arguments, and a tractable next lemma.

\textbf{Present conclusion.} There is an explicit ordinary mathematical
derivation of Gaussian boundary-to-interior estimates for a block mean
and controlled smeared channels. There is also a complete conditional
spectral-gap argument, stated and proved below. Neither establishes the
required interacting Yang-Mills estimates. No Lean-checked theorem or
successful Lean build is claimed. The original reflection-positivity
file contains unfilled assumptions and admitted proofs that require
mathematical repair before implementation.

\textbf{Requested review.} Please evaluate the Gaussian inequalities
under their literal quantifiers; verify the conditional spectral
theorem; determine whether a sufficiently strong, gauge-compatible
boundary or direct clustering hypothesis is plausible along the relevant
scaling trajectory; and identify the smallest next result that would
materially advance the construction. A useful response may reject or
replace the proposed UFBI route while preserving the abstract spectral
theorem.

The complete archive includes 13 recovered source and result files,
preserved byte for byte, new explanatory review notes, fresh numerical
logs, and this document in editable Markdown and LaTeX. The archive is a
standalone handoff, not a checkout of the upstream Lean repository. No
independent human review has yet been obtained.

\section{Evidence and claim
boundaries}\label{evidence-and-claim-boundaries}

The following distinctions govern every statement in this packet. Labels
refer to the type of evidence, not to a judgment about whether an
unproved claim is false.

\begin{longtable}[]{@{}
  >{\raggedright\arraybackslash}p{(\columnwidth - 4\tabcolsep) * \real{0.3333}}
  >{\raggedright\arraybackslash}p{(\columnwidth - 4\tabcolsep) * \real{0.3333}}
  >{\raggedright\arraybackslash}p{(\columnwidth - 4\tabcolsep) * \real{0.3333}}@{}}
\toprule\noalign{}
\begin{minipage}[b]{\linewidth}\raggedright
Component
\end{minipage} & \begin{minipage}[b]{\linewidth}\raggedright
Evidence available
\end{minipage} & \begin{minipage}[b]{\linewidth}\raggedright
What remains
\end{minipage} \\
\midrule\noalign{}
\endhead
\bottomrule\noalign{}
\endlastfoot
Gaussian estimates G1 to G8 & Ordinary derivations in the recovered
analytic note; additional AI audit & Independent mathematical review and
formalization \\
Python implementation & 12 supplied tests rerun successfully; 14
geometries and 168 channel measurements with no recorded violations &
These floating-point checks are not interval certificates or proofs of
uniformity \\
Original pilot constants & Reproduced rejection in 9 of 108 measurements
& Preserve the failed candidate as a negative control \\
Conditional spectral theorem & Complete analytic proof in this packet &
Lean implementation and assumption audit \\
Reflection-positivity scaffold & Source inspected; three admitted proof
bodies & Repair A1 and A2 statements; separate A3 dependencies \\
Upstream reconstruction interfaces & Earlier inspection recorded at a
fixed commit; selected source pages checked during packaging & No local
checkout or build; no compiled adapter \\
Wilson clustering and continuum passage & Explicit research obligations
& No proof supplied by this work \\
\end{longtable}

The Gaussian estimates concern a model with an input mass \(m>0\). They
establish consequences of that assumption. The eventual Yang-Mills task
also requires a nontrivial continuum theory with the intended field
content and axioms; the official problem statement is the target
reference {[}1{]}. A vacuum-only Hilbert space would not meet that
target, even if a spectral lower-bound statement were vacuous.

Two notation conventions prevent a silent identification of different
rates. In the Gaussian sections, \(m\) is the supplied scalar mass and
\(\gamma_*\) is a derived boundary-influence rate. In the spectral
theorem, \(\Delta\) is a common temporal autocorrelation rate. No
identity \(\Delta=\gamma_*\) is assumed. The missing transfer argument
must quantify any geometric or normalization losses.

\section{The recovered UFBI contract}\label{the-recovered-ufbi-contract}

UFBI is the working name used in the source specification {[}S{]}; no
equivalence to a standard mixing condition is assumed. For a lattice
region, a fixed central physical side \(q\), an admissible collar class
\(\mathcal C_{a,L}(M)\), and a declared measurable observation channel
\(F\), let

\[
\nu^\tau_{a,L,F}=F_*\mu^\tau_{a,L}
\]

be the channel law conditioned on collar \(\tau\). The proposed contract
is

\[
2\operatorname{TV}(\nu^\tau_{a,L,F},\nu^{\tau'}_{a,L,F})
\le \min\{2,KM e^{-\gamma D}\}.
\]

Here \(\operatorname{TV}(P,Q)=\sup_A|P(A)-Q(A)|\), so
\(2\operatorname{TV}\in[0,2]\). The constants \(K<\infty\) and
\(\gamma>0\) must be chosen before varying the admissible mesh, volume,
separation, or collars. Their permitted dependence on fixed physical
parameters and the declared observable class must be recorded. A bound
whose right side is 2 is tautological.

The original specification keeps the channel dimension bounded under
refinement. Its Gaussian supplement adds sufficient analytic control on
the whole smearing span or a trial-pairing certificate. A bounded number
of channels, bounded amplitudes of their weights, or individual energy
bounds do not by themselves establish those certificates or the density
required by the spectral theorem.

For the gauge theory, the collar class and its size parameter are still
to be defined. Compact group-valued links do not inherit the scalar
rescaling \(\tau=a\Phi\). The linear dependence on \(M\) is a Gaussian
conclusion, not an established Wilson continuity estimate.

\section{Gaussian model and explicit
estimates}\label{gaussian-model-and-explicit-estimates}

This section condenses the recovered proof {[}G{]}. The original full
derivation, numerical implementation, and tests are included in the
archive. All statements here concern the specified finite-dimensional
Dirichlet model.

\subsection{Geometry and field
normalization}\label{geometry-and-field-normalization}

Fix \(q>0\), \(m>0\), an integer \(B\ge1\), and an odd integer
\(R\ge3\). Define

\[
a=q/B,\qquad L=RB,\qquad s=(R-1)B/2,
\] \[
\Lambda_L=\{1,\ldots,L\}^4,\qquad
Q=\{s+1,\ldots,s+B\}^4,
\] \[
A=-\Delta_{\Lambda_L,D}+(ma)^2 I,\qquad \Sigma=A^{-1}.
\]

The unscaled Dirichlet Laplacian has diagonal 8 and off-diagonal \(-1\)
for interior nearest neighbors. Exterior ghost nodes are not
periodically identified. The incidence map \(\mathcal J\) sums the ghost
values adjacent to each interior node, counting every exterior link at
intersections.

Completing the square gives the interior law

\[
\mu^\tau=\mathcal N(A^{-1}\mathcal J\tau,A^{-1}).
\]

Thus the covariance is independent of the collar. Use the
four-dimensional scalar normalization \(\phi=a\Phi\) and restrict both
collars by

\[
\|\tau/a\|_\infty\le M,\qquad
\|\tau'/a\|_\infty\le M.
\]

For \(g=(\tau-\tau')/a\), put \(u=P_QA^{-1}\mathcal Jg\). The physical
mean shift in \(Q\) is \(u\), and the lattice-field shift is \(h=au\).
The exact separation from the nearest central node to an exterior ghost
node is

\[
D=a(s+1)=\frac{R-1}{2}q+a.
\]

The nominal central side \(q=aB\) differs from the ghost-to-ghost
small-box side \(a(B+1)\). The earlier exploratory code fixed the outer
box instead. Its convention \(a=1/(RB)\) is preserved in the historical
file and must not be mixed with the later fixed-\(q\) contract.

\subsection{Boundary barrier}\label{boundary-barrier}

Let

\[
\kappa_a=2\operatorname{arsinh}(ma/2),\qquad
\rho_a=\kappa_a/a,\qquad
b_a(D)=\min\{1,8e^{-\rho_aD}\}.
\]

Then \(2(\cosh\kappa_a-1)=(ma)^2\). Writing \(W\) for interior
adjacency, the convergent expansion

\[
A^{-1}=\frac1{8+(ma)^2}\sum_{n=0}^{\infty}
\left(\frac{W}{8+(ma)^2}\right)^n
\]

shows entrywise positivity. The eight face exponentials give the massive
harmonic barrier

\[
w(x)=\sum_{i=1}^4
\frac{e^{-\kappa_a x_i}+e^{-\kappa_a(L+1-x_i)}}
{1+e^{-\kappa_a(L+1)}}.
\]

It is at least 1 on each ghost face; the constant 1 is also a
supersolution. Applying the discrete maximum principle to both signs of
the extension yields

\[
\|u\|_{\infty,Q}\le\|g\|_\infty b_a(D)\le2M b_a(D).
\tag{G1}
\]

For the average \(\ell_B(v)=B^{-4}\sum_Qv(x)\), the sharper barrier is

\[
|\ell_B(u)|\le\|g\|_\infty\bar b_{B,L},
\] \[
\bar b_{B,L}=\min\left\{1,
\frac{8e^{-\rho_aD}}{1+e^{-\kappa_a(L+1)}}
\frac{1-e^{-\kappa_a B}}{B(1-e^{-\kappa_a})}\right\}.
\]

No smoothness assumption on the bounded collar values is used.

\subsection{Covariance lower bound and standardized block
mean}\label{covariance-lower-bound-and-standardized-block-mean}

For \(Y_{\rm mean}=(aB^4)^{-1}\sum_Q\phi_x\), let
\(f=B^{-4}\mathbf1_Q\), zero extended into \(\Lambda_L\). Its variance
is

\[
V_B=a^{-2}f^TA^{-1}f.
\]

For any nonzero trial vector \(v\), Cauchy-Schwarz in the \(A\)-energy
gives

\[
f^TA^{-1}f\ge\frac{(f^Tv)^2}{v^TAv}.
\]

Choose the normalized first tensor sine mode on \(Q\), zero outside, and
define

\[
\alpha_B=\frac{\pi}{2(B+1)},\quad
\lambda_B=(ma)^2+16\sin^2\alpha_B,\quad
r_B=\frac{2\cot^2\alpha_B}{B(B+1)}.
\]

The exact one-dimensional sine sums give

\[
V_B\ge v_B:=\frac{r_B^4}{a^2B^4\lambda_B}.
\]

The inequalities \(r_B\ge8/\pi^2\) and \(B^2\lambda_B\le(mq)^2+4\pi^2\)
then give the mesh-independent lower bound

\[
V_B\ge v_0:=\frac{(64/\pi^4)^2}
{q^2((mq)^2+4\pi^2)}.
\tag{G2}
\]

For completeness, on \(0\le x\le\pi/4\), convexity gives
\(\sin^2x\le2x/\pi\), while \(\sin x\le x\). At \(x=\alpha_B\) these
imply \(\cot^2\alpha_B\ge4B(B+1)/\pi^2\), which proves the stated bound
on \(r_B\).

The standardized boundary-to-mean map is

\[
\mathcal T_{\rm mean}g=
\frac{\ell_B(P_QA^{-1}\mathcal Jg)}{\sqrt{V_B}}.
\]

Consequently,

\[
\|\mathcal T_{\rm mean}\|_{\infty\to\mathbb R}
\le\frac{\bar b_{B,L}}{\sqrt{v_B}}
\le\frac{\pi^4}{64}q\sqrt{(mq)^2+4\pi^2}\,b_a(D),
\tag{G3}
\] \[
d_{\rm mean}\le
\frac{\pi^4}{4}Mq\sqrt{(mq)^2+4\pi^2}e^{-\rho_aD}.
\]

All coefficients of the scalar mean response to \(g\) are nonnegative.
Thus opposite constant collars \(\tau=+aM\), \(\tau'=-aM\) attain its
supremum over the stated amplitude class. This argument does not
establish maximization for a joint signed channel.

\subsection{Exact Gaussian distances and uniform
constants}\label{exact-gaussian-distances-and-uniform-constants}

For a linear channel with covariance \(C\) and displacement \(\delta\),
set \(d^2=\delta^TC^{-1}\delta\). For equal-covariance Gaussian laws,

\[
2\operatorname{TV}=2\operatorname{erf}\left(\frac{d}{2\sqrt2}\right)
\le\min\{2,\sqrt{2/\pi}\,d\},
\] \[
H_{\rm Hel}^2=1-e^{-d^2/8},\qquad
H_{\rm Hel}^2=1-\int\sqrt{pq}.
\]

Whitening and rotating the displacement reduces the total-variation
calculation to a one-dimensional likelihood-ratio threshold. Completing
the square gives the Hellinger identity. For singular channels, first
remove redundant coordinates, or use the inverse on their common
support. In this linear pushforward setting the displacement lies in
that support. Background conventions are recorded in {[}3{]}; the proofs
do not rely on a general unequal-covariance approximation.

Restrict \(ma\le t_*\) for fixed \(t_*>0\), and define

\[
\gamma_* =m\frac{2\operatorname{arsinh}(t_*/2)}{t_*}>0.
\]

The function \(2\operatorname{arsinh}(t/2)/t\) is decreasing, so
\(\rho_a\ge\gamma_*\). It follows that

\[
2\operatorname{TV}_{\rm mean}\le
\min\{2,K_{\rm mean}M e^{-\gamma_*D}\},
\] \[
K_{\rm mean}=\sqrt{2/\pi}\frac{\pi^4}{4}
q\sqrt{(mq)^2+4\pi^2}.
\tag{G4}
\]

The mesh rate satisfies \(\rho_a<m\) at nonzero spacing and tends to
\(m\) as \(a\to0\). A claim with exact rate \(m\) at every mesh is not
proved here. A continuum channel bound additionally requires convergence
of its means and covariances.

\subsection{Joint smeared channels and hard-edged
moments}\label{joint-smeared-channels-and-hard-edged-moments}

Use \(\langle f,g\rangle_a=a^4\sum_xf(x)g(x)\) and \(L_a=A/a^2\). For
linearly independent real weights \(\psi_i\) supported in \(Q\), define

\[
Y_i=a^3\sum_x\psi_i(x)\phi_x,\quad
C_{ij}=\langle\psi_i,L_a^{-1}\psi_j\rangle_a,\quad
\delta_i=\langle\psi_i,u\rangle_a.
\]

Suppose the whole span \(V_a\) obeys the physical energy bound

\[
\langle f,L_af\rangle_a\le\Lambda_{\rm sm}\langle f,f\rangle_a
\qquad(f\in V_a)
\]

with one \(\Lambda_{\rm sm}\) independent of mesh and outer volume. If
\(G\) is the Gram matrix of the weights, the variational inequality
gives \(C\succeq G/\Lambda_{\rm sm}\). Therefore

\[
\begin{aligned}
d_F^2&\le\Lambda_{\rm sm}\delta^TG^{-1}\delta
=\Lambda_{\rm sm}\|P_{V_a}u\|_a^2\\
&\le\Lambda_{\rm sm}q^4\|u\|_\infty^2,
\end{aligned}
\] \[
d_F\le2Mq^2\sqrt{\Lambda_{\rm sm}}\,b_a(D)
\le16Mq^2\sqrt{\Lambda_{\rm sm}}e^{-\rho_aD}.
\tag{G6}
\]

Thus \(K_F=16\sqrt{2/\pi}\,q^2\sqrt{\Lambda_{\rm sm}}\) suffices, with
rate \(\gamma_*\). There is no extra factor \(\sqrt{k}\). The factor 8
in the barrier is retained in the pure exponential coefficient.

Fixed Dirichlet sine frequencies \(1\le n_i\le J\), with \(B\ge J\),
satisfy this condition with \(\Lambda_{\rm sm}=m^2+4\pi^2J^2/q^2\). For
\(B\ge2\), the first mode and four modes with a single frequency 2 admit
the sharper value \(\Lambda_7=m^2+7\pi^2/q^2\).

The sharp block mean and coordinate dipoles have zero-extension energies
that are not uniformly bounded. The original note instead uses trial
functions \(v_j\) and matrices \(P_{ij}=\langle\psi_i,v_j\rangle_a\) and
\((E_v)_{ij}=\langle v_i,L_av_j\rangle_a\) to prove

\[
C_F\succeq P E_v^{-1}P^T.
\]

For the normalized mean plus four dipoles, with \(B\ge2\),
\(G_F\preceq I\), \(E_v\preceq\Lambda_7 I\), and

\[
s_{\min}(P)\ge s_0=32\sqrt3/\pi^4.
\]

This yields

\[
K_{\rm mean+dipoles}=
\frac{16\sqrt{2/\pi}}{s_0}q\sqrt{(mq)^2+7\pi^2}.
\tag{G8}
\]

The exact pairing sums and proof of \(s_0\) are in {[}G, section 6{]}.
For \(m=2\), \(q=1\), \(t_*=1\), the sufficient constants are:

\begin{longtable}[]{@{}lrr@{}}
\toprule\noalign{}
Channel & \(K\) & Common \(\gamma_*\) \\
\midrule\noalign{}
\endhead
\bottomrule\noalign{}
\endlastfoot
Block mean & 128.1198512355 & 1.9248473002 \\
Five sine modes & 109.1392079748 & 1.9248473002 \\
Mean plus four dipoles & 191.8093099258 & 1.9248473002 \\
\end{longtable}

These constants are conservative analytic bounds, not fitted optima. For
every bounded Borel function of a controlled joint channel,

\[
|\mathbb E_\tau f(Y)-\mathbb E_{\tau'}f(Y)|
\le2\|f\|_\infty\operatorname{TV}(\nu^\tau,\nu^{\tau'}).
\]

The displayed UFBI envelope therefore applies unchanged when
\(\|f\|_\infty\le1\). Unbounded polynomials require additional moment
control. Holding lattice collars at order one would make \(M\) grow like
\(a^{-1}\) and destroy the stated uniform constant.

\section{Conditional spectral
theorem}\label{conditional-spectral-theorem}

This is the completed analytic implication to be formalized. It assumes
a Hamiltonian and proves a spectral conclusion; it does not construct a
quantum field theory.

\subsection{Statement}\label{statement}

Let \(\mathscr H\) be a complex Hilbert space, \(H\ge0\) a densely
defined self-adjoint operator, and \(\Omega\in\operatorname{Dom}(H)\) a
unit vector with \(H\Omega=0\). Put \(\mathscr K=\Omega^\perp\). Let
\(\mathcal D\subset\mathscr K\) have dense linear span in
\(\mathscr K\).

Assume there is one \(\Delta>0\) such that, for every
\(\psi\in\mathcal D\), finite \(C_\psi\ge0\) and \(t_\psi\ge0\) exist
with

\[
0\le\langle\psi,e^{-tH}\psi\rangle
\le C_\psi e^{-\Delta t}\qquad(t\ge t_\psi).
\tag{SG}
\]

Let \(E\) be the spectral measure of \(H\), and \(P_\Omega\) the
orthogonal projection onto \(\mathbb C\Omega\). Then

\[
E([0,\Delta))=P_\Omega,\qquad
\ker H=\mathbb C\Omega,
\] \[
\sigma(H)\subseteq\{0\}\cup[\Delta,\infty),\qquad
H|_{\mathscr K}\ge\Delta I.
\]

The last inequality is understood as a quadratic-form bound. The gap is
at least \(\Delta\); an excitation at exactly \(\Delta\) is not
asserted. Nontriviality of \(\mathscr K\) is a separate requirement.

\subsection{Proof}\label{proof}

For \(\psi\in\mathcal D\), define the finite positive measure

\[
\mu_\psi(B)=\langle\psi,E(B)\psi\rangle
=\|E(B)\psi\|^2.
\]

For \(0\le r<\Delta\) and sufficiently large \(t\), positivity and the
spectral theorem give

\[
e^{-rt}\mu_\psi([0,r])
\le\int_{[0,\infty)}e^{-t\lambda}\,d\mu_\psi(\lambda)
\le C_\psi e^{-\Delta t}.
\]

Thus \(0\le\mu_\psi([0,r])\le C_\psi e^{-(\Delta-r)t}\), and letting
\(t\to\infty\) gives \(E([0,r])\psi=0\). This bounded projection
vanishes on the linear span of \(\mathcal D\) and hence on all of
\(\mathscr K\). Since \(H\Omega=0\), it is the identity on
\(\mathbb C\Omega\), so \(E([0,r])=P_\Omega\).

Choose \(r_n\uparrow\Delta\) with \(0\le r_n<\Delta\). Increasing-set
continuity of spectral projections yields \(E([0,\Delta))=P_\Omega\).
The vacuum eigenspace and its orthogonal complement reduce \(H\), and
the restricted spectral measure on \(\mathscr K\) is supported in
\([\Delta,\infty)\). This proves the form bound and spectral inclusion.
In particular \(E(\{0\})=P_\Omega\), proving vacuum uniqueness.

As a consequence, for every \(\eta\in\mathscr K\) and \(t\ge0\),

\[
0\le\langle\eta,e^{-tH}\eta\rangle
\le e^{-\Delta t}\|\eta\|^2.
\]

The proof does not need a uniform prefactor across the initial spanning
family; the positive decay rate must be common. The order of quantifiers
is \(\exists\Delta>0\;\forall\psi\;\exists C_\psi,t_\psi\).

\subsection{Observable and time
conventions}\label{observable-and-time-conventions}

With the inner product linear in the second argument, define
\(\omega(A)=\langle\Omega,A\Omega\rangle\) and
\(\psi_A=(A-\omega(A)I)\Omega\) for bounded \(A\). Then

\[
\langle\psi_A,e^{-tH}\psi_A\rangle
=\langle\Omega,A^*e^{-tH}A\Omega\rangle-|\omega(A)|^2.
\]

These positive connected temporal autocorrelations are the required
inputs. A cyclic observable algebra supplies density after centering
only if its actual vectors, including any necessary products, are
covered by the estimates. Unbounded fields need controlled domains or
directly defined reconstructed vectors.

For \(T_a=e^{-a_tH_a}\) with \(T_a\Omega_a=\Omega_a\), a decay
\(e^{-\kappa_a n}\) in transfer steps corresponds to physical energy
\(\kappa_a/a_t\). The continuum task needs a common strictly positive
physical threshold. Also,
\(\|e^{-tH}\psi\|^2=\langle\psi,e^{-2tH}\psi\rangle\): a squared-norm
rate \(\Delta\) yields energy threshold \(\Delta/2\).

The original scaffold's \texttt{HasPVMGap} only excludes the open
positive-energy window \(0<E<\Delta\). That statement alone does not
establish vacuum uniqueness. The stronger projection identity above
does.

\subsection{Exact tests of the
hypotheses}\label{exact-tests-of-the-hypotheses}

\textbf{Missing states.} On \(\ell^2(\mathbb N_0)\) take \(\Omega=e_0\),
\(He_0=0\), \(He_j=e_j\) for \(1\le j\le k\), and
\(He_{k+j}=j^{-1}e_{k+j}\) for \(j\ge1\). Every tested vector in the
first \(k\) excited coordinates has rate 1, while other excited
eigenvalues approach zero.

\textbf{No common rate.} For \(He_n=n^{-1}e_n\) and \(n\ge1\), finitely
supported centered vectors form a dense set. Each has an individual
positive exponential rate; no common positive rate exists and there is
no gap.

\textbf{Finite observation window.} For any \(H\ge0\), \(\Delta>0\), and
finite \(T\), contraction gives
\(\langle\psi,e^{-tH}\psi\rangle\le e^{\Delta T}\|\psi\|^2e^{-\Delta t}\)
on \(0\le t\le T\). A finite-time upper bound alone cannot establish the
theorem.

\section{Missing implication from boundary influence to
clustering}\label{missing-implication-from-boundary-influence-to-clustering}

The following elementary calculation exposes the needed interface. It is
added explanatory material, not a proof that the Wilson model satisfies
the interface.

Let \(F\) be a bounded interior random variable, let \(G\) be bounded
and measurable with respect to an exterior sigma algebra \(\mathcal E\),
and put \(q_F=\mathbb E[F\mid\mathcal E]\). With covariance defined as
\(\mathbb E[FG]-\mathbb E[F]\mathbb E[G]\),

\[
\operatorname{Cov}(F,G)
=\mathbb E[(q_F-c)(G-\mathbb E G)]
\]

for any constant \(c\). If \(|q_F-c|\le\epsilon\) almost surely, then

\[
|\operatorname{Cov}(F,G)|\le2\epsilon\|G\|_\infty.
\]

If a conditional channel-law bound applies to all exterior
configurations involved, its total variation controls differences of
conditional expectations. For \(F=f(Y)\) with bounded \(f\), the UFBI
convention gives
\(|q_F(\tau)-q_F(\tau')|\le\|f\|_\infty KM e^{-\gamma D}\). A fixed
reference collar supplies \(c\), provided the comparison covers the
actual conditional law.

If the bound holds only on a good event of probability at least \(1-p\)
and \(|c|\le\|F\|_\infty\), the same argument gives

\[
|\operatorname{Cov}(F,G)|
\le2\|G\|_\infty\bigl(\epsilon+2\|F\|_\infty p\bigr).
\]

The bad-collar contribution must therefore be controlled quantitatively.
The Gaussian condition \(\|\tau/a\|_\infty\le M\) is deterministic; it
is not established as a uniform typical-boundary statement for random
collars. With unrestricted scalar amplitudes and a nonzero channel
response, scaling a collar difference drives its Gaussian total
variation arbitrarily close to 1. The amplitude restriction is
substantive.

\textbf{New audit observation for review.} Suppose the collar is sampled
as unaveraged sites of a larger finite Gaussian box with this same
precision convention. Its precision \(\widetilde A\) satisfies
\(\widetilde A\preceq(16+(ma)^2)I\), so at each sampled site \(y\),

\[
\operatorname{Var}(\tau_y)=(\widetilde A^{-1})_{yy}
\ge\frac1{16+(ma)^2}.
\]

Bounding the Gaussian density by its maximum gives, even with a nonzero
mean,

\[
\Pr\{|\tau_y|\le aM\}
\le\sqrt{2/\pi}\,aM\sqrt{16+(ma)^2}\longrightarrow0.
\]

For fixed \(m,M\) the event that every raw collar site obeys the cutoff
is therefore atypical as \(a\to0\). A fixed-cutoff small-tail argument
cannot justify the desired continuum clustering for that collar law.
This deduction is added during the handoff audit, not claimed as an
earlier result or a Lean proof. It does not apply automatically to
averaged/coarse collars or compact Wilson links, and does not refute the
deterministic Gaussian theorem. It motivates asking first about an
averaged influence estimate, an appropriate boundary norm, or an
explicitly coarse boundary construction.

Before this template can supply (SG), the project must prove the
relevant Gibbs/Markov conditioning identity; specify a gauge-compatible
boundary class; control excluded collars or avoid excluding them; handle
unbounded renormalized observables; identify the reflected/adjointed OS
pairing; and relate collar separation to physical time with explicit
constants. For example, \(D\ge ct-d_0\) produces rate \(c\gamma\), not
automatically \(\gamma\).

\section{Reconstruction and limit
obligations}\label{reconstruction-and-limit-obligations}

Reflection positivity is one input to reconstruction, not the full
hypothesis set. A specific reconstruction theorem and its regularity,
covariance, symmetry, positivity, and continuity assumptions must be
identified. Osterwalder and Schrader's second paper explicitly extends
and corrects the earlier paper; its complete hypotheses should be
checked in the source rather than replaced with a shorthand slogan
{[}4{]}. Luescher's regulated transfer-matrix construction provides
relevant background, not the interacting continuum conclusion {[}5{]}.

The Gaussian boundary theorem is uniform over its declared mesh and
collar class. It does not prove convergence of the Gaussian channel
distributions, much less a Yang-Mills continuum limit. A sufficient
route for passing temporal bounds is convergence of matched, correctly
normalized correlations with finite cutoff-independent prefactors for
each fixed observable and a common physical rate. Other methods may
organize prefactors differently, but the inference must survive the
actual limits.

The recovered upstream audit records an ordered fixed-volume continuum
limit followed by a volume limit, and a distinction between unselected
correlations and a selected phase. The displayed
\texttt{WilsonRegularity.running} condition at the recorded commit
requires positive logarithmic growth of the coupling parameter as the
mesh shrinks {[}U{]}. Fixed strong-coupling lattice information alone is
not a verification of that trajectory.

If positive scalar spectral measures disintegrate as
\(\nu=\int\nu_\chi\,d\eta(\chi)\), zero subgap mass of \(\nu\) implies
zero subgap mass in almost every phase. To choose a common phase for all
relevant vectors requires a countable determining/dense family and a
proof that its phase representation is the intended one. An arbitrary
estimate on a mixed signed correlator does not transfer pointwise to
every selected phase.

The common gap theorem concerns the reconstructed Hilbert space supplied
to it, usually a vacuum sector. It does not automatically cover every
representation or superselection sector.

\section{Lean status and formalization
handoff}\label{lean-status-and-formalization-handoff}

The original \path{MassGap/ReflectionPositivity.lean} is preserved
unchanged. Direct inspection confirms three \texttt{sorry} proof bodies.
No Lean, Lake, or Elan build was run in this packaging task;
\texttt{lean} and \texttt{lake} were unavailable on the executable path.

\begin{longtable}[]{@{}
  >{\raggedright\arraybackslash}p{(\columnwidth - 4\tabcolsep) * \real{0.3333}}
  >{\raggedright\arraybackslash}p{(\columnwidth - 4\tabcolsep) * \real{0.3333}}
  >{\raggedright\arraybackslash}p{(\columnwidth - 4\tabcolsep) * \real{0.3333}}@{}}
\toprule\noalign{}
\begin{minipage}[b]{\linewidth}\raggedright
Obligation
\end{minipage} & \begin{minipage}[b]{\linewidth}\raggedright
Observed issue
\end{minipage} & \begin{minipage}[b]{\linewidth}\raggedright
Required correction
\end{minipage} \\
\midrule\noalign{}
\endhead
\bottomrule\noalign{}
\endlastfoot
A1 Wilson reflection positivity & Arbitrary reflection/support arguments
and \texttt{hGeometry\ :\ True} & Actual reflected link geometry,
admissible positive-time algebra, measure invariance, integrability, and
a Hermitian positive OS form \\
A2 transfer operator & \texttt{hQIsOSQuotient\ :\ True}; conclusion
permits the zero operator & Actual OS quotient/completion, descended
translation, bounded extension, positivity, normalization and vacuum
identities \\
A3 spectral gap & Admitted lemma hides density, spectral pairing, phase
and limit passage & Separate explicit hypotheses and proved adapters; do
not substitute Gaussian control \\
\end{longtable}

The recovered \texttt{FOUNDATION\_STATUS.md} records prior source
inspection of \path{yang-mills-group/OpenAI-YM-20261028} at commit

\path{3456e5e75cf8d4cd558955a5c7638fba2d68f464}.

It records \path{leanprover/lean4:v4.32.0} as the toolchain. This is a
historical inspection record, not a fresh successful build. The main
attachment points it identifies are \path{YM/OS/Flow.lean},
\path{YM/OS/Dynamics.lean}, \path{YM/Heat/Lift.lean},
\path{YM/Hilbert/Kolmogorov.lean}, \path{YM/Phase/Unique.lean}, and
\texttt{PublicView.lean}. The heat-lift source explicitly develops
nonnegative-energy spectral machinery without claiming a gap {[}U{]}.

\textbf{Recommended implementation order.} First prove that exponential
decay of a positive scalar Laplace transform eliminates low-energy
measure. Then prove bounded-projection annihilation on a dense span and
the vacuum projection identity. Adapt that result to the existing energy
PVM and the actual Euclidean-time pairing. Identify an unbounded
Hamiltonian and its spectrum only after the corresponding generator
relationship is established.

Acceptance requires a pinned build, exact statement review, no admitted
proof in the new bridge, and an axiom/dependency report for the bridge
and adapter. Explicit physical hypotheses may remain assumptions of a
conditional theorem, but must be displayed as such. A clean build of an
inadequate statement or a theorem depending on an assumed gap does not
satisfy this criterion.

\section{Review questions and next
work}\label{review-questions-and-next-work}

The best first reviewer is a researcher in constructive quantum field
theory or rigorous lattice gauge theory, with renormalization and
correlation-estimate experience. Functional analysis expertise is
sufficient for the abstract spectral lemma. Lean expertise is useful
after the mathematical interface is fixed.

\begin{enumerate}
\def\labelenumi{\arabic{enumi}.}
\tightlist
\item
  \textbf{Gaussian validity.} Are G1 to G8 correct under the exact
  geometry, covariance, amplitude and smearing assumptions? Are any
  stated uniformities too strong? Can the conservative bounds be
  retained as a sound control without further optimization?
\item
  \textbf{Boundary applicability.} Which gauge-invariant collar
  condition is both plausible for the Wilson measure and strong enough
  for the intended conditioning argument? Does it require
  exceptional-set estimates or a different metric?
\item
  \textbf{Temporal clustering.} What precise theorem turns that
  condition into the positive autocorrelation bound (SG), and what rate
  and prefactor losses occur?
\item
  \textbf{Density.} Which observable algebra gives a dense family of
  centered reconstructed states? Can the estimate extend from chosen
  channels to that algebra with a common rate?
\item
  \textbf{Limits and phases.} What is uniform along the actual
  coupling/mesh/volume trajectory? What passes to a common selected
  phase, and what establishes nontriviality?
\item
  \textbf{Formal fidelity.} Does the proposed PVM adapter state the
  intended conditional theorem, with the correct zero-energy sector and
  physical-time convention?
\end{enumerate}

Please return a short assessment classifying claims as correct, correct
with added hypotheses, incorrect with a counterexample, or unresolved.
For each repair, identify the precise statement or equation, the
smallest change, and any new dependency. Recommend one next analytic
lemma and one bounded formalization task, with observable completion
criteria.

\textbf{Recommended next milestone.} Complete the conditional
scalar-measure/PVM theorem and its build/axiom report, while obtaining
the Wilson-hypothesis review in parallel. Further broad Gaussian sweeps
are lower priority unless a reviewer identifies a concrete missing
check. Reconsider UFBI if the admissible collar restriction cannot
control the relevant conditional measure, or if expanding the observable
class removes the common positive rate. Direct temporal clustering can
feed the same spectral theorem.

\section{Reproducibility and fresh
checks}\label{reproducibility-and-fresh-checks}

The archive preserves the original 13 recovered files and records
SHA-256 hashes in \path{provenance/original_sources.json}. Fresh results
are separate from the originals. Source hashes embedded in the original
numerical report agree with the recovered Python files and the rerun
report.

\begin{longtable}[]{@{}
  >{\raggedright\arraybackslash}p{(\columnwidth - 2\tabcolsep) * \real{0.5000}}
  >{\raggedright\arraybackslash}p{(\columnwidth - 2\tabcolsep) * \real{0.5000}}@{}}
\toprule\noalign{}
\begin{minipage}[b]{\linewidth}\raggedright
Check performed for this handoff
\end{minipage} & \begin{minipage}[b]{\linewidth}\raggedright
Observed result
\end{minipage} \\
\midrule\noalign{}
\endhead
\bottomrule\noalign{}
\endlastfoot
Supplied unit tests & 12 passed \\
Analytic-bound sweep & 14 geometries, 56 collars, 168 channel
measurements; zero recorded violations \\
Rejected pilot replay & 108 measurements; 9 failures; strict command
returned exit status 1 as expected \\
Original implementation identity & Recorded source hashes match
recovered and executed files \\
Lean compilation and axiom audit & Not performed; compiler
unavailable \\
\end{longtable}

The pilot failures are opposite constant collars at \(R=3\) for
\(B=2,3,4\), in each of the mean, moments and full channels. Its
\(K=4\), \(\gamma=1\) candidate remains rejected. The later sufficient
constants do not retroactively validate it.

The rerun environment was Python 3.12.14, NumPy 2.3.5, and SciPy 1.17.0.
Floating-point consistency does not certify the universal analytic
quantifiers or floating-point errors. The earliest exploratory
outer-box-fixed sweep is retained but was not repeated; it is a
historical result under a different normalization.

From the extracted archive root, after installing the two packages
listed in \texttt{requirements.txt}, run:

\begin{Shaded}
\begin{Highlighting}[]
\VariableTok{OPENBLAS\_NUM\_THREADS}\OperatorTok{=}\NormalTok{1 }\ExtensionTok{python3} \AttributeTok{{-}m}\NormalTok{ unittest discover }\DataTypeTok{\textbackslash{}}
  \AttributeTok{{-}s}\NormalTok{ MassGap/UFBI }\AttributeTok{{-}p} \StringTok{\textquotesingle{}test\_*.py\textquotesingle{}} \AttributeTok{{-}v}
\VariableTok{OPENBLAS\_NUM\_THREADS}\OperatorTok{=}\NormalTok{1 }\ExtensionTok{python3}\NormalTok{ MassGap/UFBI/gaussian\_bounds.py }\DataTypeTok{\textbackslash{}}
  \AttributeTok{{-}{-}out}\NormalTok{ validation/estimate\_checks\_local.json}
\end{Highlighting}
\end{Shaded}

The README supplies the pilot replay command, its expected nonzero exit,
the file map, and instructions for rebuilding this PDF. The longer
spectral notes include additional exact counterexamples and a
version-independent formalization contract. No repository push,
publication, or message to a mathematician is part of this handoff.

\section{References and original source
map}\label{references-and-original-source-map}

\textbf{{[}G{]} Original analytic derivation.}
\path{MassGap/UFBI/GaussianBoundaryEstimate.md}, 20 November 2026.
Complete G1 to G8 proofs, units, constants, trial-pairing calculation
and numerical discussion. Original bytes preserved.

\textbf{{[}S{]} Original specification.}
\path{MassGap/UFBI/SPECIFICATION.md}, v0.1 with v0.2 analytic
supplement, 20 November 2026. The rejected pilot and later supplement
are both retained. Its early statement that scalar worst-case collars
were not known is superseded for the block mean by the later
positive-coefficient proof, as the supplement explains.

\textbf{{[}F{]} Earlier status audit.}
\path{original_notes/FOUNDATION_STATUS.md}, 20 November 2026. Records
earlier source inspection and test results; fresh checks are identified
separately above. \path{original_notes/README.md} documents the
uncompiled scaffold.

\textbf{{[}1{]} Arthur Jaffe and Edward Witten.} \emph{Quantum
Yang-Mills Theory}, official Clay problem description, especially
sections 4 and 5.
\href{https://www.claymath.org/wp-content/uploads/2022/06/yangmills.pdf}{Official
PDF}. Used to define the final target, not to justify the proposed
estimates.

\textbf{{[}2{]} Wendelin Werner and Ellen Powell.} \emph{Lecture notes
on the Gaussian Free Field}, arXiv:2004.04720v2 (2021).
\href{https://arxiv.org/abs/2004.04720v2}{Paper}. Background on discrete
and continuum Gaussian fields; its Laplacian conventions must be
translated to the diagonal-8 convention here. The explicit constants in
this packet are derived in {[}G{]}.

\textbf{{[}3{]} Luc Devroye, Abbas Mehrabian and Tommy Reddad.}
\emph{The total variation distance between high-dimensional Gaussians
with the same mean}, arXiv:1810.08693v7 (2023), section 2.
\href{https://arxiv.org/abs/1810.08693v7}{Paper}. Background for
equal-covariance TV and Hellinger identities. The paper discloses a
removed incorrect general upper bound in earlier versions; that bound is
not used here.

\textbf{{[}4{]} Konrad Osterwalder and Robert Schrader.} \emph{Axioms
for Euclidean Green's functions II}, Communications in Mathematical
Physics 42 (1975), 281-305.
\href{https://doi.org/10.1007/BF01608978}{DOI}. Publisher abstract and
bibliographic data checked; full archive retrieval was blocked during
packaging. No claim is made to have audited the complete reconstruction
hypotheses from that article in this task.

\textbf{{[}5{]} Martin Luescher.} \emph{Construction of a selfadjoint,
strictly positive transfer matrix for Euclidean lattice gauge theories},
Communications in Mathematical Physics 54 (1977), 283-292.
\href{https://doi.org/10.1007/BF01614090}{DOI};
\href{https://bib-pubdb1.desy.de/record/396349/files/7611148.pdf}{original
DESY report}. Relevant regulated positivity and transfer-matrix
construction; it is not a proof of the four-dimensional continuum mass
gap.

\textbf{{[}U{]} Upstream source references.}
\href{https://github.com/yang-mills-group/OpenAI-YM-20261028/tree/3456e5e75cf8d4cd558955a5c7638fba2d68f464}{Repository
at the recorded commit},
\href{https://github.com/yang-mills-group/OpenAI-YM-20261028/blob/3456e5e75cf8d4cd558955a5c7638fba2d68f464/ChallengeSpec.lean}{ChallengeSpec.lean},
and
\href{https://github.com/yang-mills-group/OpenAI-YM-20261028/blob/3456e5e75cf8d4cd558955a5c7638fba2d68f464/YM/Heat/Lift.lean}{YM/Heat/Lift.lean}.
The latter two source pages were checked during this handoff. The
upstream dependency graph was not rebuilt or independently certified.
The earlier audit contains the remaining pinned file links.

The separate UOR/RG contract prototype is not relied on as mathematical
evidence and is not included in this proof handoff. It would need an
explicit integration and validation step before becoming part of the
proof provenance.

\end{document}
